\chapter{Conclusion and Future Work}

This graduation project developed an AI/RAG-based live question-answer system for Gebze Technical University. The project began with the need to answer GTU-specific questions from official sources and evolved into a live broadcast assistant that can receive manual or YouTube questions, retrieve relevant GTU context, generate Turkish answers, prepare speech, and present the response through a live avatar stage.

\section{Conclusion}

The implemented system satisfies the main project objectives. It can collect and archive official GTU web and PDF sources, index documents as chunks with embeddings, retrieve evidence with a hybrid scoring method, generate institution-focused Turkish answers, keep internal answer traces, manage live question queues, and expose both admin and viewer-facing interfaces. The Docker Compose setup also provides a realistic single-server deployment path.

The evaluation demonstrates that the prototype is usable for a controlled demo. On the 12-question GTU evaluation set, the system reached an 83.33\% source hit rate and 100\% answer keyword hit rate. The top-1 source hit rate of 58.33\% and median latency of 10.07 seconds show that the system works but still needs optimization before it can be considered production-ready.

\section{Limitations}

The main limitations are:

\begin{itemize}
    \item the corpus depends on the current source profile and may miss newly published or deeply nested pages;
    \item top-1 retrieval precision is not yet high enough for fully unattended use;
    \item LLM and TTS provider calls create noticeable latency in live settings;
    \item the YouTube pipeline depends on API keys, active live chat availability, and polling intervals;
    \item the avatar animation is prototype-level and does not yet include full phoneme-level lip synchronization;
    \item there is no human approval workflow for high-stakes institutional answers.
\end{itemize}

\section{Future Work}

Future work should focus on four areas. First, retrieval can be improved with canonical URL boosts, better source-type routing, stronger document title cleaning, and more complete official GTU coverage. Second, latency can be reduced through answer caching, model selection, background precomputation for common questions, and parallel TTS preparation. Third, the live experience can be improved with richer avatar motion, better audio synchronization, improved queue transitions, and operator moderation controls. Fourth, production readiness can be strengthened with scheduled recrawling, audit logs, dashboard analytics, and human review for sensitive answers.

\section{Final Remarks}

The project shows that a university-specific live assistant can be built by combining RAG, official source archiving, modern web application design, and a broadcast-aware presentation layer. The most important design principle is not simply to make the model answer, but to make the system answer from the right evidence, at the right time, and in a form that fits the live audience.
